$p=61$: $\dim W=120$, $r=15$, $|G|=226920$, 이득 $g=63.0$, $\theta_{61}=14.754$, $\varepsilon=\tfrac{39+5\sqrt{61}}2$. unimodular 점은 120차원 짝 unimodular $SL_2(61)$-격자 $W_{61}$ 이고, 안정성은 (균형 잡힌 벡터에 대해) 최소노름 $\ge8$ 을 요구한다(노름 6 벡터는 $s$ 와 무관하게 위반; 노름 8 은 $|s|\le8$ 이어야 함).
\begin{observation}[10회차, 빌더 B]\tagN\
분리 토러스 $C_{15}$ 의 고정 직선(랭크 8), $C_5$-고정 부분격자(랭크 24), $C_3$-고정 부분격자(랭크 40)에서 $t$-원 전체(격자 간격 $0.1$)에 대해 반지름 $2\sqrt{\theta D}\cosh(0.05)$ 이하의 벡터를 전수 열거: \textbf{위반 없음}, 각각 $\min\Gamma=\theta_{61}$ 정확히(경계). 랭크 120 전체에 대해서는 $t$-격자 위의 LLL 기저와 그 쌍별 합에서 $\theta$ 미만이 없었다(약한 증거).
\end{observation}
즉 $p=37$ 을 무너뜨린 메커니즘($C_3$-고정 부분격자의 근)이 $p=61$ 에서는 나타나지 않는다. 12회차에서 부분군의 isotypic 성분($SL_2(5)$ 의 icosian 성분 등)에 대한 체계적 시험을 수행했다: 결과는 아래와 같다.
